\printedpage{19}
% head: A. Grothendieck : Dérivateurs, Ch. XII, Édition des Universités de Montpellier II et Paris VII
% foot: Version 7 décembre 2002

\section*{III\quad Remarkable classes of functors related to $W$.}

Finally, I will call \textit{``basic localizer''} a subset $W$ of Fl(Cat) satisfying the following three conditions: % [sic] appelerai -> appellerai, spelling

\begin{enumerate}
\item[(i)] Saturation W(1) (cf. reminder p. 1).
\item[(ii)] Final element axiom W(3) (page 3).
\item[(iii)] The weakened variant W(2 bis) of the localization axiom W(2) given in the corollary on page 1 (formerly it was L4$'$ in Pursuing Stacks).
\end{enumerate}

When I need more from $W$, I will say so explicitly. These three axioms have sufficed for practically everything I have done with $W$ in Pursuing Stacks, the only exception being a single use (if I recall), somewhere, of a weakened variant of W(2) (for the \textit{cartesian} $S$-functors between \textit{fibered} $S$-categories - if on every fiber it is in $W$, it is itself in $W$).

We therefore suppose henceforth given a basic localizer
\[
W \subseteq \mathrm{Fl}(\mathrm{Cat}).
\]

\textbf{Proposition 1.} \textit{Let $f : X \longrightarrow Y$ be in} \textup{Cat}. \textit{Equivalent conditions:}

\begin{enumerate}
\item[\textit{a)}] \textit{For every local isomorphism $Y' \longrightarrow Y$, $f' : X' \longrightarrow Y'$ is $\in W$.}
\item[\textit{b)}] \textit{Same condition, in the particular case $Y' = Y/y \longrightarrow Y$, so $\forall\ y \in Y$, $X/y \simeq X \times_Y (Y/y)$ is $W$-aspheric.}
\end{enumerate}

\msp{page 23}\textbf{Definition 1.} Under these equivalent conditions a), b) of proposition 1, we say that $f$ is \textit{$W$-aspheric.} We say \texttt{[that]} $f$ is \textit{$W$-coaspheric} if $f^{\circ}$ is $W$-aspheric, i.e. if the $y\backslash X$ are $W$-aspheric (for $y \in \Ob Y$).

\textsc{Proof.} It suffices to prove b) $\Longrightarrow$ a). By definition of a local isomorphism, we see that condition b) is stable under base change $Y' \longrightarrow Y$, with $Y'$ a local isomorphism (since the $Y'/y' \xrightarrow{\sim} Y/y$, hence $X'/y' \xrightarrow{\sim} X/y$). It therefore suffices to see that if $f$ satisfies b), then $f \in W$. But this is axiom W(2 bis).

\begin{center}
\begin{tikzcd}[row sep=large]
X \arrow[d] & & & \\
Y & & Y' \arrow[ll] & \\
& Y/y \arrow[ul] & & Y'/y' \arrow[ll] \arrow[ul]
\end{tikzcd}
\end{center}

\printedpage{20}
% head: A. Grothendieck : Dérivateurs, Ch. XII, Édition des Universités de Montpellier II et Paris VII
% foot: Version 7 décembre 2002

\textbf{Remark.} If $f$ has a right adjoint $g$, then for every $y \in Y$, $X/y \simeq X/g(y)$, so $X/y$ has a final object and is therefore $W$-aspheric. Thus if $f$ has a right (left) adjoint, then $f$ is $W$-aspheric ($W$-coaspheric).

But even if $f$ has no right adjoint, there is still a generalized adjoint
\[
g : Y \longrightarrow X^{\wedge}
\]
and we have
\[
X/y \simeq X/g(y) \qquad \text{(for the fully faithful inclusion $X \hookrightarrow X^{\wedge}$),}
\]
so $\mathrm{Top}(X/y)$ is identified with the induced topos $\mathrm{Top}(X)/g(y)$. The condition $f$ $W$-aspheric therefore means that the sheaves $g(y)$ in $\mathrm{Top}(X)$ are
\msp{page 24}
$W$-aspheric. For many other equivalent formulations, I refer to Pursuing Stacks.

\textbf{Proposition 2.} \textit{Let $X \xrightarrow{f} Y \xrightarrow{g} Z$ be in} \textup{Cat}. \textit{If $f$ and $g$ are $W$-aspheric ($W$-coaspheric), likewise $gf$. If $gf$ and $f$ are $W$-aspheric ($W$-coaspheric), $g$ too.}

\textbf{Recall.} A morphism $p : S' \longrightarrow S$ in Cat is a \textit{local isomorphism} if $\forall\ s' \in \Ob S'$, the induced morphism
\[
p_{s'} : S'/s' \xrightarrow{\sim} S/s \qquad (s = p(s'))
\]
is an isomorphism: $p$ induces an isomorphism on the ``localized categories'' of $S'$ and $S$ at $s'$ and $s = p(s')$ respectively. It is immediate that this notion is stable under composition, that if $gf$ and $g$ are local isomorphisms, then $f$ is one too, that every isomorphism is a local isomorphism. Moreover, local isomorphisms are \textit{stable under base change,} which follows from the fact (trivial and important) that if we have a cartesian square in Cat

\begin{center}
\begin{tikzcd}[row sep=large]
X & X' \arrow[l, "q"'] \arrow[d, "f'"] \\
S \arrow[from=u, "f"'] & S', \arrow[l, "p"']
\end{tikzcd}
\end{center}

and if $x' \in \Ob X'$ has images $x$, $s'$, $s$ in $X$, $S'$, $S$, then the diagram

\begin{center}
\begin{tikzcd}[row sep=large]
X/x \arrow[d, "f/x"'] & X'/x' \arrow[l, "q/x'"'] \arrow[d, "f'/x'"] \\
S/s & S'/s' \arrow[l, "p/s'"']
\end{tikzcd}
\end{center}

is also cartesian. (So if $p/s'$ is \texttt{[an]} isomorphism, $q/x'$ too.)

\printedpage{21}
% head: A. Grothendieck : Dérivateurs, Ch. XII, Édition des Universités de Montpellier II et Paris VII
% foot: Version 7 décembre 2002

\msp{page 25}Let us return to proposition 2. The stability of the notion of aspheric morphism under composition follows at once from form a) of the condition. Likewise for the condition that if $gf$ and $f$ are $W$-aspheric, $g$ too - we must show that if $Z' \longrightarrow Z$ is \texttt{[a]} local isomorphism, then $g'$ is in $W$, now the hypothesis implies that $g'f'$ and $f'$ are.

\begin{center}
\begin{tikzcd}[row sep=large]
X \arrow[d, "f"'] & X' \arrow[l, "r"'] \arrow[d, "f'"] \\
Y \arrow[d, "g"'] & Y' \arrow[l, "q"'] \arrow[d, "g'"] \\
Z & Z' \arrow[l, "p"']
\end{tikzcd}
\end{center}

\textbf{Corollary 1.} \textit{The inverse image of an $f : X \longrightarrow Y$ which is $W$-aspheric (resp. $W$-coaspheric) under a local isomorphism (resp. under a local co-isomorphism) is $W$-aspheric ($W$-coaspheric).}

\begin{center}
\begin{tikzcd}[row sep=huge]
X \arrow[d] & X' \arrow[l] \arrow[d] & X'' \arrow[l] \arrow[d] \\
Y & Y' \arrow[l, "\text{local iso}"] & Y'' \arrow[l, "\text{local iso}"]
\end{tikzcd}
\end{center}

Follows formally from the definition by a) and from the fact that a composite of two local isomorphisms (resp. local co-isomorphisms) is likewise.

\textbf{NB} Of course, the notion of $W$-aspheric functor is not stable under arbitrary base change. For example, the inclusion $X = \{y_0\} \longrightarrow Y$ is $W$-aspheric if and only if the discrete categories $\mathrm{Hom}(y_0, y)$ are $W$-aspheric, i.e. (provided that we have W(5), the connectedness axiom) if and only if the $\mathrm{Hom}(y_0, y)$ have cardinality 1, i.e. $y_0$ \textit{initial object} of $Y$. But if $Y \neq \{y_0\}$, taking $y$ in $Y$, $y \neq y_0$, and setting $Y' = \{y\} \longrightarrow Y$, we have $X' = X \times_Y Y' = \emptyset$ and $\emptyset \longrightarrow e$ is not aspheric!

\msp{page 26}On the other hand, we expect the $W$-aspheric morphisms to be stable under \textit{$W$-smooth} base changes (which generalize local isomorphisms), and the $W$-coaspheric morphisms under $W$-proper base changes. Cf. below.

\textbf{Corollary 2.} \textit{Let $f : X \longrightarrow Y$ in} \textup{Cat} \textit{be an} \textup{equivalence} \textit{of categories. Then $f$ is $W$-aspheric and $W$-coaspheric.}

Indeed, $f$ has a left adjoint and a right adjoint, and we use the remark on page 23.

\textbf{Proposition 3.} \textit{Let $f : X \longrightarrow Y$ be in} \textup{Cat}. \textit{Equivalent conditions:}

\printedpage{22}
% head: A. Grothendieck : Dérivateurs, Ch. XII, Édition des Universités de Montpellier II et Paris VII
% foot: Version 7 décembre 2002

\begin{enumerate}
\item[\textit{a)}] \textit{For every $y \in \Ob Y$, the inclusion functor}
\[
X_y \longrightarrow X/y \tag*{$(*)$}
\]
\textit{is $W$-}co\textit{aspheric.}
\item[\textit{b)}] \textit{For every $x_0 \in X$, the induced morphism for the colocalized categories}
\[
x_0\backslash X \longrightarrow y_0\backslash Y \qquad \textit{(where $y_0 = f(x_0)$)}
\]
\textit{has $W$-aspheric fibers.}
\item[\textit{c)}] \textit{For every morphism $\mathbf{\Delta}^1 = Y' \longrightarrow Y$, let us consider $X' = X \times_Y Y'$ over $\mathbf{\Delta}^1$, the inclusion $X'_1 \longrightarrow X'$ of the closed fiber of $X'$ (fiber over the closed point of $Y'$) is $W$-coaspheric.}
\item[\textit{d)}] \textit{For every arrow $u_0 : y_0 \longrightarrow y_1$ in $Y$, and every $x_0$ in $X_{y_0}$, the category $X(x_0;u_0)$ formed of the arrows $u' : x_0 \longrightarrow x$ in $X$ with source $x_0$, and which lift $u$} \texttt{[rather} $u_0$\texttt{]}\textit{, is $W$-aspheric.}
\end{enumerate}

\msp{page 27}\textsc{Proof.} To spell out condition a), we must spell out the category $\xi_0\backslash(X_y)$, for an arbitrary object $\xi_0 = (x_0, u_0)$ of $X/y$, $x_0 \in \Ob X$,
\[
u_0 : y_0 \longrightarrow y_1 = y \quad \text{in $Y$, with $y_0 = f(x_0)$.}
\]

This category $\xi_0\backslash X_{y_1}$ is then spelled out as the category of the $x_1$ in $X_{y_1}$, \textit{equipped} with an arrow from $\xi_0$ to the image $\xi_1 = (x_1, \mathrm{id}_{f(x_1)} = \mathrm{id}_{y_1})$ of $x_1$ in $X/y_1$ in $\xi_0$ \texttt{[``in} $\xi_0$\texttt{'' to be deleted]}, i.e. with an arrow $u$ from $x_0$ to $x_1$, such that $f(u)$ gives rise to the commutative triangle in $Y$, i.e. $f(u) = u_0$. % AUDIT resolved: regular CMTT12 editorial bracket; see corrections.csv, PDF leaf 32

\begin{center}
\begin{tikzcd}
x_1 \arrow[dd, dash] & & \\
& & x_0 \arrow[ull, "u"'] \arrow[dd, dash] \\
y_1 & y_1 \arrow[l, "\mathrm{id}"'] & \\
& & y_0 \arrow[ul, "u_0"'] \arrow[ull, "f(u)"]
\end{tikzcd}
\end{center}

So this is the category denoted $X(x_0;u_0)$ in d). We see at once that it is also isomorphic to the fiber of $x_0\backslash X \longrightarrow y_0\backslash Y$ at the element $u_0$ of $y_0\backslash Y$, referred to in b). Hence the equivalence of a) b) c). On the other hand, giving $u_0$ in $Y$ is equivalent to giving a functor $\mathbf{\Delta}^1 \longrightarrow Y$, as in c), and to express then that $\underbrace{X'_1}_{\simeq X_{y_1}} \longrightarrow X'$ is coaspheric, it suffices to express that the $x'\backslash X'_1$ ($x' \in \Ob X'$) are $W$-aspheric. Now this is trivial if $x' \in \Ob X'_1$ (initial object), so it suffices to verify it for $x'$ in $X'\backslash X'_1 = X'_0 \simeq X_{y_0}$. Then $x'$ is identified with an element $x_0$ of $X_{y_0}$, and we check at once that $x_0\backslash X'_1$ is isomorphic to the category $X(x_0;u_0)$ of d). This completes the proof.

\printedpage{23}
% head: A. Grothendieck : Dérivateurs, Ch. XII, Édition des Universités de Montpellier II et Paris VII
% foot: Version 7 décembre 2002

\msp{page 28}\textbf{Definition 2.} Under the conditions of proposition 3, we say that $f$ is \textit{$W$-proper.} We say that $f$ is \textit{$W$-smooth} if $f^{\circ}$ is $W$-proper, i.e. if the inclusion morphisms $(*)$ \texttt{[rather} $X_y \longrightarrow y\backslash X$\texttt{]} are $W$-aspheric etc. (conditions dual to those of proposition 3).

\textbf{Remark.} Condition a) of proposition 3 is satisfied in particular if the inclusion $(*)$ has a left adjoint (cf. remark on page 23). This means that the categories $X(x_0;u)$ have an initial object. This also means that the category $X$ over $Y$ is Cat-cofibered over $Y$ in the weak sense, i.e. that there exists, for $u_0 : y_0 \longrightarrow y_1$ an arrow in $Y$, and $x_0$ in $X_{y_0}$, an object $u_{0*}(x_0)$ in $X_{y_1}$ which represents the covariant functor $X_{y_1} \longrightarrow (\mathrm{Ens})$, $x_1 \longmapsto \mathrm{Hom}_u(x_0, x_1)$. Thus, if $X$ is Cat-cofibered in the weak sense over $Y$, and a fortiori if it is Cat-cofibered, then $f : X \longrightarrow Y$ is $W$-proper. Dually, if $f$ is weakly Cat-fibered, $f$ is \texttt{[}$W$\texttt{-]}smooth.

\textbf{Proposition 4.} (\footnote[43]{Say also: If $X \xrightarrow{f} Y \xrightarrow{g} Z$, if $gf$ and $f$ $W$-smooth, then $g$ is smooth at the points of $\mathrm{Im}\, f$, ditto for proper.}) \textit{The notion of $W$-proper morphism is stable under composition, under base change. Local co-isomorphisms are proper (trivial from b)), a fortiori isomorphisms are proper, closed immersions are proper.}

\textit{Dual statement for smooth morphisms: stability under base change, under composition, local isomorphisms are smooth, in particular open immersions.}

\msp{page 29}Stability under base change is immediate from form c). Stability under composition \texttt{[is]} less trivial:

\begin{center}
\begin{tikzcd}[row sep=large]
{\scriptstyle x} &[-1.5em] X \arrow[d, "f"] & x\backslash X \arrow[l] \arrow[d] & \bar{X} \arrow[l] \arrow[d] \\
{\scriptstyle y} & Y \arrow[d, "g"] & y\backslash Y \arrow[l] \arrow[d] & \bar{Y} \arrow[l] \arrow[d] \\
{\scriptstyle z} & Z & z\backslash Z \arrow[l] & \{u\} \arrow[l]
\end{tikzcd}
\end{center}

In the diagram with cartesian squares opposite \texttt{[rather whose two squares on the right are cartesian]} (deduced from $x \in \Ob X$, whence $y = f(x)$, $z = g(y) = gf(x)$, and from $u \in \Ob z\backslash Z$), we must prove that $\bar{X}$ is $W$-aspheric. We know that $\bar{Y}$ is ($g$ proper) and that $\bar{X} \longrightarrow \bar{Y}$ has $W$-aspheric fibers. It therefore suffices to use proposition 5 below, and the immediate remark % AUDIT resolved: regular CMTT12 editorial bracket; see corrections.csv, PDF leaf 33

\textbf{Lemma.} \textit{If $f$ is proper, the same holds for $x\backslash f : x\backslash X \longrightarrow y\backslash Y$, for every $x$ in $X$, $y = f(x)$.}